
et par
suite, pour tout \(k\) :
\[
  \psi^k(x+y)=\psi^k(x)+\psi^k(y).
\]

D'après (7.1.1), les \(\psi^k\) sont des polynômes isobares de poids \(k\) par rapport aux \(\lambda^i\), avec \(i=1,\ldots,k\). Donnons par exemple les valeurs de \(\psi^1\) et \(\psi^2\) :
\[
  \psi^1(x)=\lambda^1(x)=x,
  \qquad
  \psi^2(x)=x^2-2\lambda^2(x).
\]

7.2. Soit \(K\) un anneau de caractéristique \(0\)~; supposons données des applications \(\psi^k\) de \(K\) dans \(K\), additives, et telles que \(\psi^1=\operatorname{Id}_K\). On définit alors une pré-\(\lambda\)-structure sur \(K\) en posant :
\begin{equation}\tag{7.2.1}
  \lambda_t(x)=
  \exp\left(\sum_{k=1}^{\infty}(-1)^{k-1}\psi^k(x)t^k/k\right).
\end{equation}
La multiplicativité résulte de l'additivité des \(\psi^k\). Il y a donc équivalence entre la donnée d'une pré-\(\lambda\)-structure sur \(K\), et la donnée des opérations \(\psi^k\) correspondantes.

7.3. Il est clair d'après la formule (7.1.1) que les \(\psi^k\) commutent aux \(\lambda\)-homomorphismes. Si \(K,K'\) sont des pré-\(\lambda\)-anneaux de caractéristique nulle, il est nécessaire et suffisant qu'un homomorphisme \(f:K\to K'\) commute aux \(\psi^k\) pour qu'il soit un \(\lambda\)-homomorphisme~: cela résulte de 7.2, les \(\lambda^i\) s'expriment comme polynômes par rapport aux \(\psi^k\).

\textbf{Exemples 7.4.} i) Pour \(K=\mathbb Z\), muni de sa structure de \(\lambda\)-anneau, on a :
\[
  \frac{d}{dt}(\lambda_t(n))/\lambda_t(n)
  =
  n(1+t)^{n-1}/(1+t)^n
  =
  n.\sum_{k=1}(-1)^{k-1}t^{k-1}.
\]
Par suite, les \(\psi^k\) sont tous égaux à l'application identique de \(\mathbb Z\).

ii) Plaçons nous sur la catégorie \(\mathcal C_0\) des \(\mathbb Z\)-algèbres (1.8). Les 
